% META: {"id": "1SPE-SECDEG-CO-050", "chapitre": "1SPE-SECOND-DEGRE", "type_objet": "corrige", "exercice_ref": "1SPE-SECDEG-EX-050", "capacites_codes": ["C5"], "sources_inspiration": [], "status": "generated"}
% BEGIN-VERIFY
% from sympy import symbols, solveset, S, EmptySet
% x = symbols('x')
% f = x**2 + 4
% Delta = 0**2 - 4*1*4
% assert Delta == -16
% assert solveset(f, x, domain=S.Reals) == EmptySet
% END-VERIFY
\begin{corrige}{1SPE-SECDEG-EX-050}

\textbf{1. Discriminant.}

\commentaireMarge{$a = 1$, $b = 0$, $c = 4$.}

Pour l'équation $x^2 + 4 = 0$, on a $a = 1$, $b = 0$, $c = 4$, donc :
\[
\Delta = b^2 - 4ac = 0^2 - 4 \times 1 \times 4 = -16.
\]

\medskip
\textbf{2. Interprétation.}

\commentaireMarge{$\Delta < 0 \Rightarrow$ aucune racine réelle.}

$\Delta = -16 < 0$. Or, un trinôme du second degré à coefficients réels n'admet des racines réelles que si $\Delta \geqslant 0$. Donc l'équation $x^2 + 4 = 0$ n'admet \textbf{aucune solution réelle}.

\medskip
\textbf{3. Qui a raison ?}

Le \textbf{premier élève a raison}. Son raisonnement est correct : pour tout réel $x$, $x^2 \geqslant 0$, donc $x^2 + 4 \geqslant 4 > 0$. L'expression $x^2+4$ ne peut donc jamais être strictement négative, et l'inéquation $x^2 + 4 < 0$ n'a effectivement aucune solution.

Le second élève commet une erreur classique : l'affirmation « une équation du second degré a toujours deux racines » n'est vraie que dans l'ensemble des nombres \textbf{complexes} (hors programme de première), et non dans l'ensemble des réels $\mathbb{R}$. Dans $\mathbb{R}$, le nombre de racines d'un trinôme dépend du signe du discriminant : deux racines si $\Delta > 0$, une racine double si $\Delta = 0$, \textbf{aucune racine réelle si $\Delta < 0$} — ce qui est précisément le cas ici.

\end{corrige}
